% !TEX root = ../main.tex

\chapter*{出版说明}
\addcontentsline{toc}{chapter}{出版说明}

本书为《杂阿含经》（求那跋陀罗译，大正新修大藏经第 99 册）之
\textbf{研究译注本（经序重排）}：
依巴利《相应部》等平行经厘定义理，以\textbf{仿鸠摩罗什文体}之古文译出主文，
附今译意、汉译底本、平行资料与校勘说明；
阅读次序依 Anesaki 所复原、印顺《杂阿含经论会编》所采用之\textbf{卷次重排}。

\vspace{0.8em}
\noindent\textbf{版本：}
本册为系列之\textbf{V2}，与\textbf{V1《研究译注本（大正卷序）》}译文同源、体例相同，
\textbf{并行刊行、互不取代}。V1 依大正五十卷与经号 1--1362（馆藏／引用锚点）；
本册按学术卷次重排阅读，每经标明「学术序第×经 · 大正第×经」，大正经号永久保留。

\vspace{0.8em}
\noindent\textbf{本书不是}历史文献，\textbf{不声称}鸠摩罗什曾译《杂阿含》，
亦\textbf{不声称}已恢复「出土原典」之唯一经序。
「仿罗什风」指在可观察的罗什译经风格特征（删冗、四字节奏、意译圆通）下，
对阿含义理作\textbf{文体实验}；经序重排为文献学整理假说，供按相应结构阅读。

\vspace{0.8em}
\noindent 全经共 1362 经：其中 \textbf{1146} 经底本有较完整经叙；
\textbf{216} 经汉本仅为「亦如是」等一句，据巴利或同型经补最小可读文，
经题以 ◇ 标记，详见附录「重建经目录」。
正编阅读序 1359 经；第 604、640--641 经（阿育王传插入）入\textbf{附录}，不占正编序号。

\vspace{0.8em}
\noindent\textbf{制作方式：}
【正文·仿罗什风】与【今译意】依本书体例，在巴利及汉译平行经约束下，
由\textbf{大语言模型辅助生成}初稿，再经人工审读、校勘与抽检而定稿；
\textbf{非}传统意义上法师逐句手译，亦\textbf{非}未经人工核定的全自动成书。
一般读者可主读【正文】与【今译意】；研究者宜参【校勘与说明】及【平行】自行核读。

\vspace{0.8em}
\noindent\textbf{版权与署名：}
底本来自汉译佛典传统；平行、巴利及英译之来源与授权详见「文献与制作说明·版权与引用」。
引用请标明书名、版本（经序重排）、大正经号（如「第 1 经」）及所据巴利经编号（如「《相应部》22.12」）。

\vspace{0.8em}
\noindent\textbf{经序与品目：}
本册\textbf{重排}大正卷次之阅读顺序（Anesaki／印顺 48 卷流）；
大正经号不变。插入经入附录；卷二十三、二十五之暂置说明见「前言·关于经序重排」。

\clearpage
